
\documentclass[aspectratio=169]{beamer}

\mode<presentation>
\usetheme{Boadilla}
\definecolor{NTHU1}{RGB}{127,16,132}
\definecolor{NTHU2}{RGB}{228,210,231}
\definecolor{blue}{RGB}{30,90,205}
\definecolor{red}{RGB}{213,94,0}
\definecolor{green}{RGB}{0,128,0}
\definecolor{charcoalgray}{RGB}{108,104,104}
\setbeamercolor{title}{fg=NTHU1}
\setbeamercolor{frametitle}{fg=NTHU1}
\setbeamercolor{block title}{bg=NTHU1, fg=white}
\setbeamercolor{block body}{bg=NTHU2}
\setbeamercolor{structure}{fg=NTHU1}
\setbeamercolor{item projected}{fg=white}
\setbeamercolor{item}{fg=NTHU1}
\setbeamercolor{subitem}{fg=NTHU1}
\setbeamercolor{section in toc}{fg=NTHU1}
\setbeamercolor{description item}{fg=NTHU1}
\setbeamercolor{caption name}{fg=NTHU1}
\setbeamercolor{button}{bg=NTHU1, fg=white}
\setbeamercolor{caption name}{fg=NTHU1}
\usepackage{graphics}
\usepackage{tikz}
\usepackage{amsmath}
\usepackage{bbm}
\usetikzlibrary{decorations.pathreplacing}
\usepackage{geometry}
\usepackage{booktabs}
\usepackage{bookmark}
\usepackage{multirow, makecell}
\usepackage{float}
\usepackage{fancyvrb}
\usepackage{caption}
\captionsetup[figure]{singlelinecheck = false, format= hang, justification=centering}
\captionsetup[subfigure]{singlelinecheck = false, format= hang, justification=raggedright}
\usepackage{subcaption}
\usepackage{adjustbox}
\usepackage{hyperref}
\usepackage{threeparttable}
\usepackage{appendixnumberbeamer} 
\usepackage[T1]{fontenc}
\usepackage[default]{lato} 
\usepackage{smartdiagram}
\smartdiagramset{
  back arrow disabled = true,
  module minimum width=1cm,
  module x sep = 3,
  uniform arrow color=true,
}
\usepackage{xcolor}
\usepackage{colortbl}
	\definecolor{light-gray}{gray}{0.99}

\newenvironment{wideitemize}{\itemize\addtolength{\itemsep}{10pt}}{\enditemize}
\newenvironment{wideenumerate}{\enumerate\addtolength{\itemsep}{10pt}}{\endenumerate}
\newenvironment{widedescription}{\description\addtolength{\itemsep}{10pt}}{\enddescription}
\hypersetup{
colorlinks=true,
linkcolor=NTHU1,
filecolor=green, 
urlcolor=blue,
}
\beamertemplatenavigationsymbolsempty
\setbeamercolor{author in head/foot}{bg=white, fg=NTHU1}
\setbeamercolor{title in head/foot}{bg=white, fg=NTHU1}
\setbeamercolor{date in head/foot}{bg=white, fg=NTHU1}
\setbeamercolor{section in head/foot}{bg=white, fg=NTHU1}
\setbeamercolor{headline}{bg=white}
\setbeamertemplate{footline}{
    \leavevmode%
    \hbox{%
        \begin{beamercolorbox}[wd=.333333\paperwidth,ht=2.25ex,dp=1ex,center]{date in head/foot}%
            \usebeamerfont{date in head/foot}%\insertdate
        \end{beamercolorbox}%
        \begin{beamercolorbox}[wd=.333333\paperwidth,ht=2.25ex,dp=1ex,center]{title in head/foot}%
            \usebeamerfont{title in head/foot}\insertshorttitle
        \end{beamercolorbox}%
        \begin{beamercolorbox}[wd=.333333\paperwidth,ht=2.25ex,dp=1ex,right]{date in head/foot}%
            \usebeamerfont{date in head/foot}\insertframenumber{}\hspace*{2em}
        \end{beamercolorbox}}%
        \vskip0pt%
    }
    %% May need to script the introduction!!!!!!
%\setbeamercolor{page number in head/foot}{fg=NTHU1}
\setbeamertemplate{section in toc}[sections numbered]
\setbeamertemplate{subsection in toc}{\leavevmode\leftskip=3em\rlap{\hskip-1.75em\inserttocsectionnumber.\inserttocsubsectionnumber}
\inserttocsubsection\par}
\setbeamerfont{subsection in toc}{size=\footnotesize}
%\setbeamertemplate{headline}{%
  %\begin{beamercolorbox}[ht=5.5ex]{section in head/foot}
    %\vskip2pt\insertnavigation{0.33\paperwidth}\vskip2pt
  %\end{beamercolorbox}%
%}
\newenvironment{transitionframe}{\setbeamercolor{background canvas}{bg=white}\setbeamertemplate{footline}{} \begin{frame}[noframenumbering]}{\end{frame}}


\makeatletter
\let\@@magyar@captionfix\relax
\makeatother


\title[]{Conflict and Peace} % Change this regularly
\subtitle[]{Week 2: Economic Incentives behind Conflicts}
\author[Lee]{Seung-hun Lee }
\institute[]{Taipei School of Economics\\ National Tsing Hua University}

\date[]{September 12th, 2025}

\begin{document}
\begin{transitionframe}
\titlepage
\end{transitionframe}


%%% Color slides for section headers: Use for colloquium version (The ones bewteen \iffals and \fi)

\section{Overview of the topic}
\begin{frame}
\frametitle{Stylized facts about conflicts (Blattman, Miguel, 2010)} 
Many conflicts take form of internal conflicts
\begin{itemize}\setlength\itemsep{3pt}
\item Roughly 20-30\% of countries have active civil war/conflicts during 1960-2000s
\item At a given year: Kill approx 270,000 ppl \& 8.44 mil disability-adjusted life-years $\Downarrow$
\item Slows down GDP growth (2\% yr) \&  human capital, infrastructure, and institutions $\Downarrow$
\item Disproportionately affect poorer regions 
\end{itemize}\medskip
Economic roots of (civil) wars
\begin{itemize}\setlength\itemsep{3pt}
\item Poverty and weak economic growth correlated with conflict onset
\item Resource abundance raises risk of conflicts via higher rents to be seized
\item Individuals more inclined to join conflict when returns to work and peace are low 
\item Weak state capacity to contain conflict leads to longer duration
\end{itemize}
\end{frame}

\begin{frame}
  \frametitle{Share of wars of different types}
  \begin{figure}
    \includegraphics[keepaspectratio,width=0.7\textwidth]{fig5.png}
    \caption*{Source: Ray, Esteban (2017) ``Conflict and Development", \textit{Annual Reviews of Economics},9,263-293}
  \end{figure}
  
\end{frame}



\begin{frame}
  \frametitle{Share of countries with active conflicts each year}
  \begin{figure}
    \includegraphics[keepaspectratio,width=0.8\textwidth]{bm2010_f1.jpeg}
  \end{figure}
\end{frame}

\begin{frame}
  \frametitle{Disproportionate distribution of conflicts}
  \begin{figure}
    \includegraphics[keepaspectratio,width=0.8\textwidth]{bm2010_f2.jpeg}
  \end{figure}
\end{frame}


\begin{frame}
  \frametitle{Why do people fight? (Blattman, 2022)}
War is costly, yet people fight due to 
\begin{itemize}\setlength\itemsep{3pt}
\item \textbf{Unchecked interests:} Decisionmakers that benefit from wars not held accountable
\item \textbf{Intangible interests:} Preference for values such as identity-based vengeance and glory
\item \textbf{Uncertainty:} Lack of information on others' (and own) benefits and constraints
\item \textbf{Commitment Problem:} Inability to trust promises to constrain power over time
\item \textbf{Misperceptions:} Mistaken beliefs that lead to errors in judgement
\end{itemize} 
These are not mutually exclusive \\
\bigskip
In all, economic calculuus explains why conflicts occur\\
\bigskip
We explore how economics and political science approaches the causes of conflict with these baseline reasons in mind
\end{frame}

\begin{frame}
  \frametitle{Why? Still an ongoing problem}
\begin{figure}
  \centering
  % First row
  \begin{subfigure}{0.4\textwidth}
    \includegraphics[width=0.75\linewidth, keepaspectratio]{fig1.png}
    \caption{Israel-Palestine War}
  \end{subfigure}
  \begin{subfigure}{0.4\textwidth}
    \includegraphics[width=0.75\linewidth]{fig2.png}
    \caption{Russia-Ukraine War}
  \end{subfigure}
  
  % Second row
  \begin{subfigure}{0.4\textwidth}
    \includegraphics[width=0.75\linewidth]{fig3.png}
    \caption{Religious conflict in Burkina Faso}
  \end{subfigure}
  \begin{subfigure}{0.4\textwidth}
    \includegraphics[width=0.75\linewidth]{fig4.png}
    \caption{Martial law attempt (South Korea)}
  \end{subfigure}

\end{figure}
\end{frame}


\begin{transitionframe}
\frametitle{Roadmap for today}
\tableofcontents
\end{transitionframe}

\section{Resource and commodity shocks affect conflicts}



\subsection{Caselli, Morelli, Rohner (2015): The Geography of Interstate Resource Wars}
\begin{transitionframe}
\begin{center}
  \begin{huge}
    \textcolor{NTHU1}{%
  Caselli, Morelli, Rohner (2015): \\ The Geography of Interstate Resource Wars
    }
  \end{huge}
\end{center}
\end{transitionframe}


\begin{frame}
\frametitle{Q: Resource endowment \& Location $\Rightarrow$ Interstate conflict} 
Natural resources as triggers of wars
\begin{itemize}\setlength\itemsep{3pt}
\item Many individual cases studies, but lack systematic causal estimation
\item Theoretical predictions
\begin{itemize}\setlength\itemsep{3pt}
\item Resource endowments closer to the border $\Rightarrow$ $\Uparrow$ Interstate Conflicts
\item The farther the resources, more costlier to claim them
\end{itemize}
\item Empirical analysis
\begin{itemize}\setlength\itemsep{3pt}
\item Negative correlations between distance and conflicts documented
\item Possible omitted variable bias: Resource-endowed countries innately belligerent?  
\end{itemize}
\end{itemize}
Key innovation and takeaways
\begin{itemize}\setlength\itemsep{3pt}
  \item Incorporate likelihood of conflicts over geographical setting with asymmetrical payoff
  \item Model how motives of countries respond to resource values and Acquisition costs
\end{itemize}
\end{frame}

\begin{frame}
\frametitle{Correlation between Conflict and distance to resources} 
\begin{figure}
  \centering
  \includegraphics[width=0.7\textwidth]{cmr(2015).png}
\end{figure}
\end{frame}

\begin{frame}
\frametitle{Setting up theoretical framework} 
Setup for the Model (payoffs in the matrix)
\begin{center}
\begin{tabular}{|l|c |c |}
\hline
& Country B: Peace  &Country B: Conflict\\ \hline
 Country A: Peace   & $0,0$ & $x+c_A, -x + c_B$ \\ \hline
 Country A: Conflict &$x+c_A, -x + c_B$ & $x+c_A, -x + c_B$\\ \hline
\end{tabular}
\end{center}
\begin{itemize}\setlength\itemsep{3pt}
\item $x$ is the payoff of conflict to $A$, $-x$ is the loss from conflict to $B$  (asymmetry)
\begin{itemize}
\item If $x>0$: $A$ is the expected winner
\end{itemize}
\item $c$ is the net benefit (benefit $-$ cost) of conflict
\begin{itemize}
\item Usually $c<0$: Conflicts are very costly..!
\end{itemize}
\item Peace persists if
\begin{itemize}
  \item for A: $0\geq x+c_A $ and for B: $0\geq -x + c_B \Rrightarrow c_B \leq x \leq -c_A$
  \item i.e. Conflicts are costly ($c_A$ and $c_B$ far below zero) and gains from conflicts ($x$) is small
\end{itemize}
\end{itemize}
\end{frame}

\begin{frame}
\frametitle{Takeaways from the framework} 
Testible predictions
\begin{itemize}\setlength\itemsep{3pt}
\item Higher asymmetric payoff ($x\Uparrow $) $\Rrightarrow$ Conflict rises
\begin{itemize}
\item One country has oil > None have oil
\item Two countries have oil > None have oil
\item One country has oil > Two countries have oil (smaller asymmetry in the latter)
\end{itemize} \medskip
\item Conflict inversely related with cost of acquisition (distance) 
\begin{itemize}
\item One country has oil: Conflict rises when oil is close
\item Both have oil: Changes in distance that increses oil presence asymmetry increase conflicts
\end{itemize} \medskip
\item Overall: Payoff asymmetry and lesser cost of acquisition makes conflicts likely
\end{itemize}
\end{frame}

\begin{frame}
\frametitle{Data and Identification} 
Sources of data
\begin{itemize}\setlength\itemsep{3pt}
\item Dyadic Militarized Interstate Dispute dataset with 0-5 scale for conflict intensity

\item Explanatory variables
\begin{itemize}
\item Georeferenced oil data (PETRODATA) $+$ Distance from borders from CShapes dataset
\item Control: land area, GDP per capita, population, democracy, trade, among others
\end{itemize}
\end{itemize} \medskip
Regression Approach (by country pair ($d$)-year ($t$))
\[
\begin{aligned}
Y_{d,t+1}&=\alpha+\beta\text{One}_{d,t}+\gamma{\text{One}\times\text{Dist}}_{d,t} \\ 
&+\delta\text{Both}_{d,t}+\eta{\text{Both}\times\text{MinDist}}_{d,t}+{\text{Both}\times\text{MaxDist}}_{d,t}+\xi X_{d,t}+u_{d,t}
\end{aligned}
\]
\begin{itemize}\setlength\itemsep{3pt}
\item (None/One/Both) countries in pair $d$ has resources
\item Control for distance from border (MinDist: Smaller of the distances where Both $=1$)
\item Comparison relative to country pairs without any oils (None)
 \end{itemize}
\end{frame}

\begin{frame}
\frametitle{Key Results} 
\begin{figure}
\includegraphics[width=0.8\textwidth]{cmr2015_t1.jpeg}
\end{figure}
\begin{itemize}
\item When all are controlled, Oil presence $\to$ $\Uparrow$ Conflcit and $\Uparrow$ Dist $\to$ $\Downarrow$ Conflict
\end{itemize}
\end{frame}

\begin{frame}
\frametitle{Mechanism and discussion} 
Findings are robust to
\begin{itemize}\setlength\itemsep{3pt}
\item Different measures of conflicts (Full scale war only, disupte intensity)
  \item Controlling for oil production/reserve amounts
\item Estimation methods (Logit, Probit regressions)
\item Reducing sample to cases where there were no border changes
\end{itemize} \medskip
Key takeaways from the paper
\begin{itemize}\setlength\itemsep{3pt}
\item Asymmetrical oil presence $\to$ $\Uparrow$ Conflicts
\item Distance matters: $\Uparrow$ Distance to oil $\to$ $\Downarrow$ Conflict
\item Most dangerous: Only one country has oil, and it is close 
\item Thus, conflicts increase when gains are asymmetric and costs of acquisition are low
\end{itemize}
\end{frame}



\subsection{Dube, Vargas (2013): Comodity Price Shocks and Civil Conflict: Evidence from Colombia}
\begin{transitionframe}
\begin{center}
  \begin{huge}
    \textcolor{NTHU1}{%
  Dube, Vargas (2013): \\ Comodity Price Shocks and Civil Conflict: Evidence from Colombia
    }
  \end{huge}
\end{center}
\end{transitionframe}

\begin{frame}
\frametitle{Q: How do different income shocks affect conflicts?} 
Relationship between income and conflict: Complicated
\begin{itemize}\setlength\itemsep{3pt}
\item Reduce conflicts: Wages $\Uparrow$ \& Labor supplied to conflict $\Downarrow$ (opporunity cost effect)
 \item Increase conflcits: Spoils to fight over $\Uparrow$ (rapacity channel)
 \begin{itemize}
\item Presumably, different types of income shocks lead to different effect on conflicts
\end{itemize}
\end{itemize}
\medskip
This paper leverages different sources of income shocks to answer this dilemma
\begin{itemize}\setlength\itemsep{3pt}
  \item Colombia: Long history of civil conflicts, with high escalation in 1990s
\item Agricultural goods (labor-intensivity $\Uparrow$) vs Natural resource prices (labor-intensivity $\Downarrow$)
 \begin{itemize}
\item Coffee vs Oil: Colombia's biggest exports
\end{itemize}
\item Use variation in municipal conflict incidence and commodity distribution in Colombia
\end{itemize}
\end{frame}


\begin{frame}
\frametitle{History of civil conflicts in Colombia} 
Long history of civil wars
\begin{itemize}\setlength\itemsep{3pt}
\item Left-wing guerrilla (FARC, ELN) vs govt vs right-wing paramilitaries (Oficina Cartel etc)
\item Rooted in \textit{La Violencia} in 1948, current form took place in 1960s
\item Heightened in 1990s, ongoing low-intensity fights despite peace agreement in 2016
\end{itemize}
\medskip
What did they fight over?
\begin{itemize}\setlength\itemsep{3pt}
\item Appropriate resource returns, government revenues through various channels
\item Coffee returns huge: One of the largest exporters in the world
\item Actively recruited rural workers by paying wages over minimum wage
\end{itemize}
\end{frame}


\begin{frame}
\frametitle{Data and empirical strategy} 
Sources of (some of the) data
\begin{itemize}\setlength\itemsep{3pt}
\item Conflict Analysis Resource Center (CERAC) for civil war episodes: 1988-2005
\item Commodity distribution: Agricultural land area, oil/gold production, coca production
\item Price: International Coffee Organization and Internatinal Financial Statistics  
\end{itemize} \medskip
Empirical strategy 
\begin{itemize}\setlength\itemsep{3pt}
\item DID with IV for coffee prices:  Colombian conflicts $\to$ prices directly (reverse causality)
\[
Y_{jrt} = \alpha_j + \beta_t + \delta_r t + \gamma\text{Coca}_{jr}t+\lambda(\text{Oil}_{jr}\times PO_t)+\rho(\widehat{\text{Coffee}_{jr}\times PC_t}) + \phi X_{jrt}+e_{jrt}
\]
\item Rainfall $\times$ Temperature $\times$ Exports of other big producers as IV
\begin{footnotesize}
\[
\text{Coffee}_{jr}\times PC_t = \alpha_j + \beta_t + \delta_r t + \gamma\text{Coca}_{jr}t+\lambda(\text{Oil}_{jr}\times PO_t) + \sum_{m=0}^1\sum_{n=0}^1\theta_{mn}(R_{jr}^m\times T_{jr}^n\times E_{t})+ \phi X_{jrt}+e_{jrt}
\]
\end{footnotesize}
\end{itemize}
\end{frame}

\begin{frame}
\frametitle{Descriptive findings} 
\begin{figure}
\centering
\begin{subfigure}{0.49\textwidth}
  \includegraphics[keepaspectratio, width=\textwidth]{dv2013_f4.jpeg}
  \caption{Coffee price \& trend of violence in coffee vs non-coffee municipalities}
\end{subfigure}
\begin{subfigure}{0.49\textwidth}
  \includegraphics[keepaspectratio, width=\textwidth]{dv2013_f5.jpeg}
  \caption{Oil price \& trend of violence in oil vs non-oil municipalities}
\end{subfigure}
\end{figure}
\end{frame}

\begin{frame}
  \frametitle{Overall effect on violence}
  \begin{figure}
  \includegraphics[keepaspectratio, width=0.9\textwidth]{dv2013_t2.jpeg}
  \end{figure}
\end{frame}

\begin{frame}
  \frametitle{Leading mechanism behind the outcomes}
  \begin{figure}
  \includegraphics[keepaspectratio, width=0.8\textwidth]{dv2013_t3.jpeg}
  \end{figure}
\end{frame}

\begin{frame}
\frametitle{Discussion} 
Robustness of the outcomes and ruling out other mechanisms
\begin{itemize}\setlength\itemsep{3pt}
\item Results consistent when substitute commodities are used
\begin{itemize}\setlength\itemsep{3pt}
\item Oil $\to$ Coal and Gold
\end{itemize}
\item Not driven by changes in population composition following conflict-induced migration
\item Not affected by decreasing government presence or enforcement (increases instead)
\item Results unaffected when coca-production is accounted for
\end{itemize}
\medskip
Key takeaways from the paper
\begin{itemize}\setlength\itemsep{3pt}
\item Different commodity price shocks trigger different channels of conflict
\item Labor-intensive price $\Uparrow$ $\to$ Opportunity cost of forgoing labor $\Uparrow$ $\to$ Conflict $\Downarrow$
\item Less labor-intensive price $\Uparrow$ $\to$ Maximizable rents $\Uparrow$ $\to$ Conflict $\Uparrow$

\end{itemize}
\end{frame}

\subsection{McGuirk, Burke (2020): The Economic Origins of Conflict in Africa}
\begin{transitionframe}
\begin{center}
  \begin{huge}
    \textcolor{NTHU1}{%
  McGuirk, Burke (2020): \\ The Economic Origins of Conflict in Africa
    }
  \end{huge}
\end{center}
\end{transitionframe}

\begin{frame}
\frametitle{Q: How do price shocks differently affect producers \& consumers?}
Explaining how economic shocks can affect conflict is tricky
\begin{itemize}\setlength\itemsep{3pt}
\item Shocks to opportunity cost also affect spoils and capacity to repel insurgents
\item Income may affect conflict, and vice versa (reverse causality)
\item Both income and conflict may be affected by omitted factors (legal institutions)
\end{itemize}\medskip
This paper differentiates incidence of conflicts by actors and motives 
\begin{itemize}\setlength\itemsep{3pt}
\item Leverage food price shocks and differences in food-production across households
\begin{itemize}\setlength\itemsep{3pt}
  \item Grain price $\Uparrow$ $\to$ Income $\Uparrow$ for producers ($\because$ Higher returns on food production)
  \item Grain price $\Uparrow$ $\to$ Real Income $\Downarrow$ for consuming households ($\because$ Spend more on food)
\end{itemize}
\item Conflict may take different forms
\begin{itemize}\setlength\itemsep{3pt}
  \item Conflict over land, workers, and resources (factor conflict) - aim to displace population
  \item Conflict over splitting spoils (output conflict) - more transient
\end{itemize}
\end{itemize}\medskip
\end{frame}

\begin{frame}
\frametitle{Theoretical framework}
Setting
\begin{itemize}\setlength\itemsep{3pt}
\item Consider an economy of producers (landownders) and consumers (work or fight)
\item Food price changes affect returns from land \& relative wages between work vs fighting
\end{itemize}\medskip
Factor conflict
\begin{itemize}\setlength\itemsep{3pt}
\item Producer price $\Downarrow\to$ Farming returns $\Downarrow$. Consumper prices $\Uparrow\to$ Income $\Downarrow$ \& Fight $\Uparrow$
\item Crop price $\Uparrow$: Reduce factor conflict for producers ($\because$ choose farming over attacking)
\item  Crop price $\Uparrow$: Increase conflict for consumers ($\because$ choose high-paying armed groups)
\end{itemize}\medskip
Output conflict
\begin{itemize}\setlength\itemsep{3pt}
\item Both consumers and producers can appropriate returns from food price increase
\item With higher spoils available,  Crop price $\Uparrow\to$ Conflict $\Uparrow$ for both
\end{itemize}\medskip
\end{frame}

\begin{frame}
\frametitle{Data and Empirical Strategy}
Data sources at the grid level data
\begin{itemize}\setlength\itemsep{3pt}
\item Conflict: UCDP, ACLED, and Afrobarometer
\begin{itemize}\setlength\itemsep{3pt}
\item Factor conflict: involve more organized groups using UCDP
  \item Output conflict: riots and violence against civilians using ACLED and Afrobarometer
\end{itemize}
\item Food/consumer price: Spatial distribution of crops + annual world price (UNFAO)
\end{itemize}\medskip
Empirical strategy leverage price and conflict variations at cell $i$, country $l$, and year $t$
\[
\begin{aligned}
\text{Factor conflict}_{ilt} &=\alpha_i + \sum_{k=0}^2 \beta^p_{ilt-k}\text{PPI}_{ilt-k}+\sum_{k=0}^2 \beta^m_{ilt-k}\text{CPI}_{ilt-k} + \gamma_l\times t + e_{ilt}\\
\text{Output conflict}_{ilt} &= \alpha_i+ \sum_{k=0}^2 \theta^p_{ilt-k}\text{PPI}_{ilt-k}+\sum_{k=0}^2 \theta^m_{ilt-k}\text{CPI}_{ilt-k} + \gamma_l\times t + e_{ilt}
\end{aligned}
\]
\begin{itemize}\setlength\itemsep{3pt}
\item Hypotheses: $\beta^p<0$, but $\beta^m, \theta^p, \theta^m>0$
\end{itemize}\medskip
\end{frame}

\begin{frame}
\frametitle{Descriptive trends in conflict incidence}
\begin{figure}
\centering
\begin{subfigure}{0.49\textwidth}
  \includegraphics[keepaspectratio, width=\textwidth]{mb2020_f1a.jpeg}
  \caption{Factor conflicts and food price}
\end{subfigure}
\begin{subfigure}{0.49\textwidth}
  \includegraphics[keepaspectratio, width=\textwidth]{mb2020_fb2.jpeg}
  \caption{Output conflicts and food price}
\end{subfigure}
\end{figure}
\end{frame}


\begin{frame}
\frametitle{Regression results}
\begin{figure}
\centering
\begin{subfigure}{0.425\textwidth}
  \includegraphics[keepaspectratio, width=0.9\textwidth]{mb2020_t1.jpeg}
  \caption{Factor conflicts and food price}
\end{subfigure}
\begin{subfigure}{0.425\textwidth}
  \includegraphics[keepaspectratio, width=0.9\textwidth]{mb2020_t2.jpeg}
  \caption{Output conflicts and food price}
\end{subfigure}
\end{figure}
\end{frame}


\begin{frame}
\frametitle{Notable heterogeneity of effects}
Effects are weaker in more economically developed areas (factor conflicts)
\begin{itemize}\setlength\itemsep{3pt}
  \item Uses nightlight luminosity to proxy economic development (Table 3)
  \item Coefficient on CPI $\times$ luminosity $<0$
  \item Incidence of conflicts after price shocks less likely in more developed areas
  \item These areas suffer less income shocks compared to poorer areas
\end{itemize} \medskip
Shocks to food prices have greater impact than cash crop prices
\begin{itemize}\setlength\itemsep{3pt}
  \item Cash crops: Non-food crops that does not impact real income of consumers
  \item Food crop price $\Uparrow\to$  Conflict $\Uparrow$, Cash crop price $\Uparrow\to$  Conflict $\Downarrow$ 
\end{itemize} \medskip
These results suggest household-level economic shocks determine decision to fight
\end{frame}

\begin{frame}
\frametitle{Discussion}
Income shocks matter
\begin{itemize}\setlength\itemsep{3pt}
  \item Uses price changes that generate opposing income shocks within countries 
  \item Finds opposing trends in the incidence of conflicts within countries
  \item Rules out the explanation that state's capacity to deter insurgency explains conflicts
  \item Advances the discussion using much more localized data than in the past
\end{itemize} \medskip
Where to go forward
\begin{itemize}\setlength\itemsep{3pt}
  \item Need to locally tailor policy response aimed at minimizing violence
  \begin{itemize}\setlength\itemsep{3pt}
    \item Workfare programs for rural areas (keep opportunity cost of joining conflicts high)
    \item Buffer stocks or insurance products to shelter consumers from negative income shocks
    \end{itemize} 
  \item Political/historical grievances may still play a role in conflict incidence (Next week)
  
\end{itemize} \medskip
\end{frame}



\subsection{Cao, Chen (2022): Rebel on the Canal: Disrupted Trade Access and Social Conflict in China, 1650-1911 [student presentation]}
\begin{transitionframe}
\begin{center}
  \begin{huge}
    \textcolor{NTHU1}{%
  Cao, Chen (2022): \\ Rebel on the Canal: Disrupted Trade Access and Social Conflict in China, 1650-1911 \\  (student presentation)
    }
  \end{huge}
\end{center}
\end{transitionframe}





\section{Economic opportunities to individuals explain conflicts}
\subsection{Shapiro, Vanden Eynde (2023): Fiscal Incentives for Conflict: Evidence from India's Red Corridor}
\begin{transitionframe}
\begin{center}
  \begin{huge}
    \textcolor{NTHU1}{%
 Shapiro, Vanden Eynde (2023): \\ Fiscal Incentives for Conflict: Evidence from India's Red Corridor }
  \end{huge}
\end{center}
\end{transitionframe}

\begin{frame}
\frametitle{Q: How does tax revenue structures affect conflict outcomes?} 
Fiscal revenue structures can shape incentives and prizes for conflict
\begin{itemize}\setlength\itemsep{3pt}
\item Differentiated revenue and expenditure stream in decentralized contexts
\item Resource royalties are shared with small number of local governments
\item Distribution of such royalties can lead to violence, $\to$ Govt fail to internalize benefits
\end{itemize} \medskip
This paper
\begin{itemize}\setlength\itemsep{3pt}
\item Utilizes variations in resource royalties across time and types of resources
\begin{itemize}\setlength\itemsep{3pt}
\item 10-fold increase in iron royalty to states 
\item Compare with other resources witout such hikes
\end{itemize}
\item Document rise in state-vs-insurgent conflicts, with aim to take control of these zones
\end{itemize}
\end{frame}

\begin{frame}
\frametitle{Maoist Conflict in India} 
Fiscal structure in India
\begin{itemize}\setlength\itemsep{3pt}
\item Central govt sets royalty rates on minerals
\item State govt owns resources and receives revenues, but cannot set own rates
\item Iron ore royalty to states $\times10$ since `09: Royalty $\Uparrow$ and Price $\Uparrow$ 
\end{itemize}\medskip
Mining in India
\begin{itemize}\setlength\itemsep{3pt}
\item Largest producers of iron, bauxite, and coal
\item Maoist rebels benefit: Take control of mining areas + Extortion from miners
\end{itemize}\medskip
Maoist conflicts (Naxalite)
\begin{itemize}\setlength\itemsep{3pt}
\item Communist uprising against that started as low intensity in 1960s
\item Intensified in 2004 following merger of two major factions
\item Claimed 5,000+ lives in `07-`13, prevalent in center-east regions (Red Corridor region)
\end{itemize}
\end{frame}

\begin{frame}
\frametitle{Red Corridor in India} 
\begin{figure}
  \begin{subfigure}{0.49\textwidth}
  \includegraphics[keepaspectratio,width=\textwidth]{sve2023a.jpeg}
  \end{subfigure}
  \begin{subfigure}{0.49\textwidth}
  \includegraphics[keepaspectratio,width=\textwidth]{sve2023_a1.jpeg}
  \end{subfigure}
\end{figure}
\end{frame}

\begin{frame}
\frametitle{Data and Empirics} 
Data Sources
\begin{itemize}\setlength\itemsep{3pt}
\item Conflict data: South Asia Terrorism Portal (2017-13)
\begin{itemize}\setlength\itemsep{3pt}
\item Maoist $\to$ Police, and vice versa, at district level
\end{itemize}
  \item Mineral endowments: Geological Survey of India + satellite imagery
\end{itemize}\medskip
Empirical Strategy: Diff-in-diff using introduction of royalty reforms
\begin{itemize}\setlength\itemsep{3pt}
\item OLS with Fixed effects: Compare districts with different iron endowments
\[
y_{ist}=\alpha_i + \beta \text{Iron}_{is}\times\text{Post}_t+ \gamma \text{Iron}_{is}\times\log{\text{Iron price}}_t + \mu_{st}+e_{ist}
\]
\item Main parameter: $\beta$ (Effect of reforms on iron-producing areas)
\item Control for international iron price: Value of iron itself could affect conflicts 
\end{itemize}
\end{frame}


\begin{frame}
\frametitle{Incidence and intensity of violence rises} 
\begin{figure}
  \includegraphics[keepaspectratio,width=0.9\textwidth]{sve2023_t2.jpeg}
\end{figure}
\begin{itemize}\setlength\itemsep{3pt}
\item Mean of $\log(\text{Iron})$: 1.7 $\to$ What happens by moving from no-iron to mean-iron area?
\item Incidence: Police-on-Maoist violence 12 p.p (=0.068 $\times$ 1.7)$\Uparrow$
\item Intensity: Increase by 12\% and 10\% each (=0.056 $\times$ 1.7)
\end{itemize}
\end{frame}

\begin{frame}
\frametitle{Red Corridor in India} 
\begin{figure}
  \includegraphics[keepaspectratio,width=0.9\textwidth]{sve2023_t3.jpeg}
\end{figure}
\begin{itemize}\setlength\itemsep{3pt}
\item Increase in illegal mining by 15 p.p
\item State's counterinsurgency $\uparrow$ for rents or miners help Maoists to keep illegal activities
\end{itemize}
\end{frame}

\begin{frame}
\frametitle{Discussion} 
What drives violence (or not)
\begin{itemize}\setlength\itemsep{3pt}
\item Effect fades out: State seeems to have established control 
\item Not driven by resource prices: No statistical relationship with conflicts
\item This effect not shown with other types of resources
\item Changes in incentives: Reflected in illegal mining activities
\end{itemize}
Takeaways
\begin{itemize}\setlength\itemsep{3pt}
\item Fiscal revenue structure can explain conflicts
\item Driven by increased security efforts or miners funding conflict to avoid royalties
\item There are cases where bandits take over, especially in state absence \\ (See Sanchez de la Sierra 2020, to be covered later in semester)
\end{itemize}

\end{frame}


\subsection{Blattman, Annan (2016): Can Employment Reduce Lawlessness and Rebellion?}
\begin{transitionframe}
\begin{center}
  \begin{huge}
    \textcolor{NTHU1}{%
 Blattman, Annan (2016): \\ Can Employment Reduce Lawlessness and Rebellion? }
  \end{huge}
\end{center}
\end{transitionframe}

\begin{frame}
\frametitle{Q: Does incentives to individuals affect conflict participation} 
Many fragile states provide work to deter interest in conflicts and illicit activities
\begin{itemize}\setlength\itemsep{3pt}
\item Aim to stimulate employment by providing capital and training 
\item Decrease incentives for illegal work
\item Provide income to integrate soldiers back into economy
\end{itemize} \medskip
However, ex-fighters (High-risk men) are different
\begin{itemize}\setlength\itemsep{3pt}
\item Comparative advantage in violence over human/social capital
\item Generally show lack of interest in existing rehabilitation programs
\end{itemize}\medskip
This paper experimentally assesses the effect of providing capital on conflict participation
\begin{itemize}\setlength\itemsep{3pt}
\item Experiments effects of providing agricultural training and capital inputs to ex-soldiers
\item Highlights that long-run economic incentives can deter conflict participation
\end{itemize}\medskip
\end{frame}

\begin{frame}
\frametitle{Liberian Civil wars} 
History
\begin{itemize}\setlength\itemsep{3pt}
\item Civil wars in 1989-1996 and 1999-2003
\item Killed 10\% of population, with many young men recruited to fight
\item By 2008: Ex-soldiers engaged in illicit activities while clustered in remote `hotspots'
\end{itemize} \medskip
Experiment
\begin{itemize}\setlength\itemsep{3pt}
\item Takes place in 2009 in two regions (Bong and Sinoe)
\item Trains agricultural skills and social skills during program
\item After graduation, participants provided with farming tools worth around \$125 
\item Note that some were given cash (implementation mistakes)
\end{itemize}\medskip
\end{frame}


\begin{frame}
\frametitle{Research Design} 
\begin{figure}
\includegraphics[keepaspectratio,width=0.45\textwidth]{ba2016_f1.jpeg}
\end{figure}
\end{frame}

\begin{frame}
\frametitle{Empirical Strategy} 
\begin{itemize}
\item Program assignment randomized: OLS regression between treated vs control group
\item Randomized within blocks defined by training site, rank, community
\end{itemize}
\begin{figure}
\includegraphics[keepaspectratio,width=0.5\textwidth]{ba2016_t1.jpeg}
\end{figure}
\end{frame}


\begin{frame}
\frametitle{Shift away from illicit activities to agriculture} 
\begin{figure}
\includegraphics[keepaspectratio,width=0.67\textwidth]{ba2016_t3.jpeg}
\end{figure}
\begin{itemize}
\item More interest in agriculuture, but not complete shift from illicit activities
\end{itemize}
\end{frame}

\begin{frame}
\frametitle{Decreased interst in violence following the program} 
\begin{figure}
\includegraphics[keepaspectratio,width=0.67\textwidth]{ba2016_t4.jpeg}
\end{figure}

\end{frame}

\begin{frame}
\frametitle{Other notable outcomes and discussions} 
Other findings
\begin{itemize}\setlength\itemsep{3pt}
\item Limited effect on networking and integration
\begin{itemize}\setlength\itemsep{3pt}
\item Combintation of economic interventions and other social interventions required?
\end{itemize}
\item Impacts of training alone without capital lower
\begin{itemize}\setlength\itemsep{3pt}
\item Provision of \textit{long-run} economic incentives required for effective rehabilitation
\end{itemize}
\end{itemize} \medskip
Takeaway
\begin{itemize}\setlength\itemsep{3pt}
\item Individual's decision to participate in illicit activities respond to economic margins
\item However, these should come with long-run horizons and may not affect integration
\item Economic and behavioral treatment suggested for effective rehabiliation \\ (also generalizable to developed countries)
\end{itemize}\medskip
\end{frame}

\subsection{Berman et al (2011): Do Working Men Rebel? Insurgency and Unemployment in Afghanistan, Iraq, and the Philippines [student presentation]}
\begin{transitionframe}
\begin{center}
  \begin{huge}
    \textcolor{NTHU1}{%
  Berman et al (2011): \\ Do Working Men Rebel? Insurgency and Unemployment in Afghanistan, Iraq, and the Philippines [student presentation]}
  \end{huge}
\end{center}
\end{transitionframe}

\begin{frame}
\frametitle{Takeaways and remaining questions} 
Economic incentives do determine decision to fight
\begin{itemize}\setlength\itemsep{3pt}
  \item Resource endowments shape returns from participating in war vs peaceful activities
  \item High opportunity costs deter participating in conflicts
  \item Individual incentives matter: Largely comes from benefits of labor market vs conflicts
\end{itemize} \medskip
Unanswered questions
\begin{itemize}\setlength\itemsep{3pt}
  \item How to design incentives to deter interest in conflicts?
  \item Some institutions prevent incentives from spilling into conflicts: What? How?
  \item Other determinants of conflict: Cultural factors, misperceptions from misinformation, among other things
\end{itemize}
\end{frame}


%%%%%%%%%%%
\end{document}
